\section{Proofs for the reduction}\label{app:hardness}
This appendix contains the proofs of the lemmas of Section~\ref{sec:hardness} that are not given there. The computer-assisted steps (CA1)--(CA4) are marked as in Section~\ref{sec:hardcomp}.

\subsection{Proof of the chain lemma}
\begin{proof}[Proof of Lemma~\ref{lem:hchain}]
\emph{(a) Each $P_a$ is a ZE partition.} This is a conjunction of local conditions: every member, and every hard pair that is an edge, involves vertices of at most three consecutive cells, and the pairs of vertices that must be apart are only constrained through clusters, which have vertices in at most three consecutive cells. A condition at a cell therefore depends only on the relation of the position $a$ to this cell (smaller by at least $3$, or $a-j\in\{2,1,0,-1,-2\}$, or larger), on the distance to the left end up to a constant, and on the distance to the right end up to a constant. Windows at the left end do not depend on $m$, windows at the right end depend only on the distance to the right end, and interior windows are isomorphic. All window types occur for $m\le 20$. \textbf{(CA1)}: for every $7\le m\le60$ and every $a$ the partition $P_a$ was checked by computer to have cost exactly $|H|$ with one edit per member and no hard pair edited (no failure). By the window argument this covers all $m\ge7$.

\emph{(b) There are no other ZE partitions.} \textbf{(CA2)}. Process the cells $0,1,\dots,m$ from left to right. After cell $c$, a \emph{partial ZE partition} is a partition of the vertices of the cells $\le c$ in which every hard pair and every member with all its vertices in these cells is respected. A vertex of cell $c'$ can be together with a vertex of a cell $c''\ge c'+3$ only if the pair is soft or a hard edge, which never happens (soft pairs join cells at distance at most $2$ and the hard edges join vertices in cells at distance at most $1$). So only the partition induced on the vertices of the last two cells matters for the future, together with the information which blocks already contain an older vertex (such a block is \emph{closed}). The state after cell $c$ is the set partition of the vertices of the cells $c-1,c$ with its set of closed blocks, with blocks renamed canonically and vertices described by (type, cell offset). At every interior cell there are exactly six reachable states, and the transition matrix between consecutive interior cells is the same matrix
\[
M=\begin{pmatrix}0&0&0&0&1&1\\0&0&0&0&0&0\\1&1&1&1&0&0\\0&0&0&0&0&0\\0&0&0&0&1&1\\0&0&0&0&0&0\end{pmatrix},\qquad\chi_M(x)=x^4(x-1)^2 ,
\]
with a Jordan block of size two at the eigenvalue $1$ (computed with exact integer arithmetic; the equality of the matrices was verified at the cells $7,\dots,19$ of $m=26$, and the transition rule at a cell involves only members and hard pairs within three consecutive cells, while the member list is translation invariant in the interior, so the matrix is the same at every interior cell of every $m$). The number $N_m$ of ZE partitions is $N_m=u^{\top}M^{\,m-c_0}v_m$ where $u$ is the vector after the left cap, which does not depend on $m$, and $v_m$ is determined by the right cap and depends only on the distance to the right end, hence is the same for all $m$ after the shift by the exponent. Since $x^4$ divides $\chi_M$, the generalised eigenspace of the eigenvalue $0$ is annihilated by $M^4$, and $M^{t}$ is for $t\ge4$ an affine function of $t$ (the Jordan block at $1$). Hence $N_m$ is an affine function of $m$ for $m\ge m_1$. A direct computation gives $N_m=m$ for $7\le m\le40$, so $N_m=m$ for all $m\ge7$. Together with (a), which provides $m$ distinct ZE partitions (their spine words $R^aML^{m-1-a}$ differ), these are all of them.
\end{proof}
\begin{remark}
As a consistency check that is not used in the proof, an independent SAT enumeration of all ZE partitions of the thin chain returned exactly the $m$ partitions $P_a$ for $m=7,\dots,14$, for $m=18,26,42$, and for a sample of larger $m$.
\end{remark}

\runinhead{Translation invariance, precisely.} The member list of $\mathrm{OH}_m$ consists of a \emph{left cap} (the members on the first few cells), a \emph{right cap} (the members on the last few cells) and the interior families $P_j,G_j,F_j$. The map $s_t\mapsto s_{t+2}$, $e_t\mapsto e_{t+1}$ shifts every cell by one and is an isomorphism between the sub-instances induced on the cells $\{j,j+1,j+2\}$ and $\{j+1,j+2,j+3\}$ whenever both windows avoid the caps. Consequently: (i) the constraints at the cells of the left cap involve only the cap members and the first interior members, the same for every $m$, so \emph{left cap windows are independent of $m$}; (ii) the constraints near the right end involve the right cap and the last interior members, the same for all $m$ after the shift of the index by $m$, so \emph{right cap windows depend only on the distance to the right end}; (iii) every other window is interior and isomorphic to every other interior window. In (CA1) every window type of (i)--(iii) occurs for $m\le20$, so a check of all $m\in[7,60]$ is a check of every window type, and the set of window types does not grow for larger $m$. In (CA2) the vector $u$ after the left cap is independent of $m$, the interior transfer matrix is the same $M$ at every interior cell (by the isomorphism above), and the right-cap vector $v$ depends only on the distance to the right end, so $N_m=u^{\top}M^{\,m-c_0}v$ with $u,v,c_0$ fixed. The finite computation and the affine form of $M^t$ then determine $N_m$ for all $m\ge7$. This is a finite certificate combined with the invariance argument above, not a closed-form proof.
\subsection{Proof of the profile lemma}
\begin{proof}[Proof of Lemma~\ref{lem:hprof}]
Read off the definition of $P_a$: the vertex $s_{4u+2}$ is $s_{2j+2}$ for $j=2u$ and $s_{2j}$ for $j=2u+1$. If $a<2u$ the cluster $C_{2u}$ has the form $\{s_{4u+1},s_{4u+2}\}\cup\{\varepsilon(2u)\}$ and $C_{2u+1}\not\ni s_{4u+2}$; if $a=2u$ the cluster is $C_{2u}=\{s_{4u},s_{4u+1},s_{4u+2},\varepsilon(2u)\}$; if $a=2u+1$ it is $C_{2u+1}=\{s_{4u+2},s_{4u+3},s_{4u+4},\varepsilon(2u+1)\}$; if $a\ge2u+2$ it is $C_{2u+1}=\{s_{4u+2},s_{4u+3},\varepsilon(2u+1)\}$. At the caps, the label $\varepsilon$ takes the exceptional values stated before Lemma~\ref{lem:hchain}. \textbf{(CA3)}: the statement was in addition verified by computer for all even $m$ with $8\le m\le60$ and all $u$ ($1{,}836$ cases), including the groups $u=0$ and $u=U-1$.
\end{proof}
\subsection{Proof of the parameter lemma}
\begin{proof}[Proof of Lemma~\ref{lem:hpar}]
(a) Every pair between guests and super-vertices, and every pair inside a super-vertex that is used as a pad or host pair, lies in exactly one member because different guests use different classes of the one-factorisation at the same $v$ and a pair of a matching lies in exactly one pad (or host) member; Latin members are pair-disjoint as shown above and have no pair inside a super-vertex; the single-vertex members $(x,y,z)$, $(z,y,z')$ use pairs between guests and anchors only. Every member $(a,b,c)$ has the edges $ab,bc$ and the non-edge $ac$ by definition, and these adjacencies are consistent because a soft pair belongs to one member.

(b) \textbf{(CA4)} A hand count from Lemma~\ref{lem:hprof}: the sets $W$ contain $h_u$, $\Pi_0,\Pi_3$ (pad), $\Pi_0$ (X), $\Pi_3$ (Y) or all $\Pi_i$ (All) of a single group $u$, so a super-vertex is used only by guests of the group $u$, except $s_{4u+4}=s_{4(u+1)}$, which occurs in $\Pi_2$ of the group $u$ and $\Pi_1$ of the group $u+1$. A group carries: a tile group the permanent guest $z$ and a lock guest $y$ (both $W\ni h_u$); a separator group (with $W_z=\mathrm{Pad}_u$) the permanent guest and at most two comparator guests (one with the chain as first chain, using $X_u$, one as second chain, using $Y_u$); any other group only $z$. At the host $h_u$: the permanent guest $z$ (host matching) plus either the lock guest $y$ of a tile group (a tile group carries no comparator), giving $u=2$, or, at a separator group, the up to two comparator guests (one with $h_u$ in $X_u$ of the first chain, one in $Y_u$ of the second), giving $u\le3$; a group is a tile group, a separator group or neither, never two of these. At $s_{4u+1},\varepsilon^-_u$: $z$ plus $y$ or the comparator guest of the first chain: at most $2$; at $s_{4u+3},\varepsilon^+_u$: $z$ plus $y$ or the comparator of the second chain: at most $2$; at $s_{4u},s_{4u+4}$: only in $\mathrm{All}$, i.e.\ only the permanent guests of tile groups, at most one from each of two groups: at most $2$. If $k=1$ there are no comparators, and if $k=2$ each chain has at most one neighbour, so the host has at most $2$. So a host carries $2$ (host matching and lock-guest pad) or $3$ (host matching and two comparator pads at a separator) matchings, and every non-host super-vertex at most $2$. The value $3$ occurs, at a separator host, only for $k\ge3$; the count of $u_{\max}$ by computer (thin-guest model, all tiles present, $1\le k\le5$, $2\le n\le4$) gave $2$ for $k\le2$ and $3$ for $k\ge3$, never more. The thin member counts: a vertex lying in four thin members is an $e$-vertex, and the hosts $s_{4u+2}$ lie in three; this, and the above counts, were confirmed by computer for $k\le4$ and $n\le3$ (maximum of $u(v)$ and of the number of thin members through every vertex, with the exceptional values at the caps).

(c) A vertex of $K_v$ has glue degree at least $p-1-u(v)$. For $p\ge8$ and $u(v)\le3$ this is at least $p-4\ge(p-1)/2$, hence the glue graph is connected (diameter at most $2$). For $p=6$ we have $k=1$, so $u(v)\le2$ and the degree is at least $3\ge(p-1)/2$. At most $3\le p-1$ matchings are used at any vertex.

(d) A thick vertex $(v,i)$ lies in $p$ Latin members of each thin member through $v$, and in $u(v)$ guest members. By (b) this is at most $4p+2$ at a vertex in four thin members and at most $3p+3\le4p+2$ otherwise. The sizes: $2k^2$ chains with $O(n^2)$ cells each and $O(p)$ thick vertices per thin vertex; each thin member gives $p^2$ Latin members; guests contribute $O(p)$ members each.
\end{proof}
\subsection{Proofs of the pathwidth and pair-maximality lemmas}
\begin{proof}[Proof of Lemma~\ref{lem:hpw}]
Members of different chains share a vertex only through a comparator guest or through its pad members (they share vertices with Latin members of both chains); tunnels are hard pairs and are not members. This gives the components. Define a \emph{level} $\lambda(M)$ of a member: for a Latin member of a thin member, the smallest cell of the thin member (all three cells lie in $[\lambda,\lambda+2]$); for every member belonging to a guest of the group $u$ (host, pad, lock, comparator members), $\lambda=2u$. A guest of group $u$ has pad and host vertices in cells $[2u,2u+2]$ (the vertices $s_{4u},\dots,s_{4u+4},\varepsilon^\pm_u$), so for every thick vertex $(v,i)$ all members through it have levels in $[\mathrm{cell}(v)-2,\mathrm{cell}(v)]$, and all members through a guest have the same level. Hence two members that share a vertex have levels differing by at most $2$. For a component $Q$ put $B_j=\{M\in Q:\ j\le\lambda(M)\le j+2\}$; every member lies in the bags $j\in[\lambda(M)-2,\lambda(M)]$, an interval, and two members sharing a vertex have a common bag; so $(B_j)_j$ is a path decomposition. At one level a chain has at most $6$ thin members (at most $6$ at the right cap, at most $3$ elsewhere), hence at most $6p^2$ Latin members. At an even level $2u$, the guest members of one chain are at most $3.5p$ (permanent guest) $+\,2.5p+1$ (lock) in a tile group and $2.5p$ in a separator group, and the comparator guests of a row or column at a separator group contribute at most $k-1$ guests with at most $3p+1$ members each. So a level of a component has at most $6kp^2+k(6p+1)$ members, and $|B_j|\le3\cdot6kp^2+2k(6p+1)\le24kp^2$ for $p\ge6$. The bound does not depend on $n$ because the number of members at one level is a function of $p$ and $k$ only: the cell structure is the same in every chain, guests and comparators of a group sit at the same level $2u$ in every chain, and the level range of a member is at most $2$.
\end{proof}
\begin{remark}
As a consistency check that is not used in the proof, the path decomposition with bags $\{M:\lambda'(M)\le t\le\lambda''(M)\}$ (where $[\lambda'(M),\lambda''(M)]$ is the interval spanned by the coordinates of the vertices of $M$) was verified by computer for small parameters: for every vertex the intervals of the members through it have a common point, every member spans at most $2$ cells, and for $k=1$, $p=8$ and $n=2,\dots,5$ the largest bag has between $1296$ and $1304$ members, independent of $n$, as the proof states.
\end{remark}
\begin{proof}[Proof of Lemma~\ref{lem:hpm}]
An induced $P_3$ with three hard pairs consists of two hard edges sharing a vertex and a hard non-edge. The hard edges form the graph $\Gamma_h$ whose components are: (a) for a thin vertex without a thin hard edge, the glue of $K_v$; (b) for each of the four thin hard edges $uv$ of a chain, $K_u\cup K_v$ with the glue inside both and all edges between them; (c) the tunnels $xx'$ (the anchors carry no other hard edge, guests carry none). In the components (a) and (b) the pairs that are not adjacent are exactly the pad and host pairs inside a super-vertex, which are soft; a component (c) is a single edge. So in every component every non-adjacent pair is soft, and an induced $P_3$ with two hard edges has its three vertices in one component of $\Gamma_h$, hence its third pair is soft or an edge. Hence no induced $P_3$ has all pairs hard. (The statement was also confirmed by computer for small parameters.)
\end{proof}
